% Chương 4 -- hệ thống và API. Căn cứ: code/API as-built branch dev.
\section{Hệ thống và giao diện lập trình ứng dụng}\label{sec:system-api}

Kết quả dự báo chỉ có giá trị thực hành khi có thể đi vào một quy trình xử lý rõ ràng. Hệ thống của đề tài vì vậy được tổ chức thành hai luồng tách biệt. Luồng ngoại tuyến nạp dữ liệu Olist, biến đổi kho dữ liệu, huấn luyện và chấm điểm theo lô. Luồng trực tuyến nhận điểm đã được ghi, tạo hàng đợi để nhân viên chăm sóc khách hàng xử lý, và lưu vết thao tác. API không tự nạp dữ liệu, huấn luyện hoặc chấm điểm khi khởi động; các công việc này phải được gọi tường minh. Sự tách biệt này giúp tránh việc một yêu cầu web làm thay đổi ngầm dữ liệu hoặc phiên bản mô hình.

\subsection{Kiến trúc và ranh giới dữ liệu}

Hệ thống gồm PostgreSQL, dịch vụ FastAPI, ứng dụng Angular và các job Python, được đóng gói bằng Docker Compose. PostgreSQL giữ ba vùng dữ liệu có vai trò khác nhau: Bronze/Silver/Gold phục vụ kho dữ liệu; schema \texttt{ml} giữ release, snapshot dự đoán và ngưỡng; schema \texttt{app} giữ người dùng, case, chu kỳ xử lý, liên lạc, quyết định, tham số và audit. Angular gọi REST/JSON đến FastAPI; FastAPI chỉ điều phối request, xác thực, phân quyền và ánh xạ lỗi, còn luật nghiệp vụ nằm ở domain/use case và adapter PostgreSQL thực hiện giao dịch.

Điểm nối giữa hai luồng là snapshot dự đoán bất biến. Job chấm điểm tạo điểm ở mức item, sau đó tổng hợp điểm đơn bằng luật \texttt{max-v1}. Snapshot mang phiên bản model, phiên bản ngưỡng, thời điểm chấm điểm và nguồn dữ liệu. Workflow có thể thay đổi người được giao, trạng thái, liên lạc hoặc quyết định; không có API sửa trực tiếp score hay risk level. Cách tách này cho phép sau này đối chiếu kết quả xử lý với đúng điểm đã tạo hàng đợi, thay vì đối chiếu với một điểm đã bị thay đổi bởi thao tác vận hành.

\subsection{Xác thực, vai trò và vòng đời case}

API sử dụng Bearer JWT; mật khẩu được băm và token có thời hạn. Hai vai trò là \texttt{SPECIALIST} và \texttt{MANAGER}. Specialist xem hàng đợi, tự nhận đơn mức thấp/trung bình, ghi lần liên lạc và tải minh chứng. Manager phân công, mở lại case, đặt quyết định, xem dashboard/audit và thay đổi tham số vận hành. Kiểm tra quyền được thực hiện ở backend, không chỉ dựa vào việc ẩn nút trên giao diện.

Case có bốn trạng thái: \texttt{PENDING}, \texttt{IN\_PROGRESS}, \texttt{COMPLETED} và \texttt{SKIPPED}. Specialist chỉ nhận case đang chờ, chưa có người nhận, không phải mức HIGH và còn năng lực xử lý. Case HIGH phải được Manager phân công. Khi hoàn tất hoặc bỏ qua, case đóng; Manager có thể mở lại với lý do, khi đó một chu kỳ mới với SLA mới được tạo nhưng lịch sử cũ vẫn được giữ. Giao dịch nhận đơn kiểm tra xung đột và năng lực trong backend để hai người cùng thao tác không thể cùng sở hữu một case.

\subsection{Thiết kế API}

Mọi endpoint nghiệp vụ dùng tiền tố \texttt{/api/v1}; dữ liệu JSON dùng \texttt{snake\_case}, thời điểm dùng ISO~8601 có timezone. Lỗi nghiệp vụ được ánh xạ thành HTTP 403/404/409/422 với mã lỗi để giao diện phân biệt lỗi quyền, trạng thái, dữ liệu đầu vào và xung đột. Bảng~\ref{tab:api-groups} tóm tắt các nhóm API chính; bảng chỉ mô tả contract phục vụ luận văn, không thay thế đặc tả DTO chi tiết của dự án.

\begin{table}[htbp]
\centering\small
\caption{Các nhóm endpoint chính của hệ thống}
\label{tab:api-groups}
\begin{tabular}{@{}p{3.0cm}p{5.1cm}p{5.0cm}@{}}
\toprule
Nhóm & Endpoint tiêu biểu & Chức năng và quyền \\
\midrule
Xác thực và hồ sơ & \texttt{POST /auth/login}, \texttt{GET/PATCH /auth/me}, \texttt{POST /auth/logout} & Phát hành JWT, lấy/cập nhật hồ sơ và kết thúc phiên ở client. \\
Hàng đợi & \texttt{GET /orders}, \texttt{/orders/summary}, \texttt{/orders/filter-options}, \texttt{POST /orders/\{id\}/claim} & Lọc, phân trang, xếp ưu tiên; Specialist tự nhận case hợp lệ. \\
Workflow case & \texttt{GET /orders/\{id\}}, \texttt{POST /orders/\{id\}/assign}, \texttt{/complete}, \texttt{/skip}, \texttt{/reopen} & Đọc chi tiết, phân công, đóng hoặc mở lại; Manager có quyền rộng hơn. \\
Liên lạc và quyết định & \texttt{POST /files}, \texttt{/contacts}, \texttt{/comments}, \texttt{/decisions}, \texttt{/model-notes} & Ghi liên lạc, minh chứng, trao đổi và quyết định; các bản ghi là append-only. \\
Giám sát và cấu hình & \texttt{GET /dashboard}, \texttt{/audit-events}, \texttt{GET/POST /parameters}, \texttt{GET /config} & Dashboard, audit và phiên bản tham số; các thao tác quản trị yêu cầu Manager. \\
\bottomrule
\end{tabular}
\end{table}

Với hàng đợi, thứ tự mặc định là điểm giảm dần, sau đó hạn SLA gần nhất và thời gian chờ. API hỗ trợ lọc theo mức rủi ro, trạng thái, người xử lý, người bán, danh mục, khoảng score, SLA và nhu cầu quyết định. Điều này biến score thành một đầu vào cho quyết định phân bổ công việc, thay vì tự động gắn nhãn hay tự động liên hệ khách hàng.

\subsection{Dấu vết vận hành và giới hạn hiện tại}

Mỗi thao tác thay đổi workflow sinh audit event; contact, comment, decision và parameter version không bị ghi đè. Upload minh chứng được kiểm tra kiểu/kích thước và gắn với lần liên lạc, thay vì chấp nhận URL tùy ý. Cấu trúc này cung cấp dữ liệu cần thiết để kiểm tra SLA, nguồn gốc quyết định và về sau đánh giá liệu hàng đợi có giúp phát hiện vấn đề tốt hơn hay không.

Phạm vi hiện tại là hệ thống hỗ trợ quyết định cục bộ/demo, không phải dịch vụ dự báo production. Bản demo dùng snapshot mang trạng thái \texttt{SYNTHETIC\_DEMO}; điểm hiển thị không được trình bày là suy luận xác nhận từ mô hình đã chọn. Ngoài ra, một số điểm còn cần hoàn thiện trước khi công bố rộng: metric test niêm phong chưa có, chưa có SHAP, chế độ \texttt{BUSINESS\_HOURS} mới được lưu cấu hình nhưng SLA backend vẫn tính theo giờ liên tục, và hợp đồng giữa một số màn hình/frontend với endpoint cần tiếp tục kiểm thử tích hợp. Các giới hạn này là lý do phần đánh giá hệ thống chỉ xác nhận luồng nghiệp vụ/API và phân biệt nó với đánh giá chất lượng mô hình ở mục~\ref{sec:analysis}.
